\documentclass[12pt,a4paper]{article}

\usepackage[utf8]{inputenc}
\usepackage[T1]{fontenc}
\usepackage[margin=2.5cm]{geometry}
\usepackage{hyperref}
\usepackage{microtype}
\usepackage{parskip}
\usepackage{lmodern}

\hypersetup{colorlinks=true, urlcolor=blue!60!black}

\pagestyle{empty}

\begin{document}

\begin{flushright}
Kundai Farai Sachikonye \\
Auf der Lichtung 19, 82178 Puchheim \\
\href{https://github.com/fullscreen-triangle}{github.com/fullscreen-triangle} \\
\texttt{kundai.sachikonye@bitspark.com} \\
(+49) 1626386344 \\[6pt]
16 May 2026
\end{flushright}

\vspace{1em}

Prof.\ Roland Rad \\
Institute of Molecular Oncology and Functional Genomics \\
TUM Klinikum Rechts der Isar \\
Munich, Germany

\vspace{1.5em}

\textbf{Re: PhD / Postdoc Position --- Cancer Biology and/or Bioinformatics (26\_05\_001)}

\vspace{1em}

Dear Prof.\ Rad,

I am applying for a position in the Institute of Molecular Oncology and Functional
Genomics, specifically in the area of computational large-scale data integration.
My background is computational biology with doctoral research at TUM, and the tools
I have built for large-scale genomic data analysis are publicly deployed.
The sequence homology search tool built on the
\href{https://github.com/fullscreen-triangle/gospel}{Gospel} framework is available at
\href{https://vivid-symbolism.vercel.app/search}{vivid-symbolism.vercel.app/search}
and can be used directly without installation.

\textbf{Large-scale genomic data integration.}
Gospel is a whole-genome variant analysis and RNA-seq integration framework designed
for the scale that massively parallel CRISPR and transposon screens generate.
It processes million-scale VCF files at $10^6$ variants per second using Bayesian
network inference, Gaussian mixture models, and variational Bayes for multi-omics
integration. It has been validated against 1000 Genomes Phase~3, gnomAD v3.1.2,
and UK Biobank, and handles the full pipeline from variant calling and RNA-seq
quantification through pathway analysis (KEGG, Reactome) and functional annotation.
The multi-dimensional integration problem that arises when combining CRISPR screen
readouts, somatic mutation profiles, transcriptomics, and clinical outcomes is
precisely what the framework was designed for.

\textbf{Cancer biology and immune mechanisms.}
My manuscript \textit{On the Thermodynamic Consequences of Categorical Completion
on Biological Systems} develops a model of adaptive immune selectivity that predicts
cancer-associated proteins as primary targets from structural first principles.
MHC binding affinity correlates with the categorical richness of the parent
protein's isoform and post-translational modification landscape
($r = 0.73$, $p < 10^{-12}$, $n = 2{,}847$ IEDB peptides); the model achieves
AUC 0.81 against NetMHCpan's 0.68, and explicitly predicts collagen isoforms and
tumour-associated antigens as primary immune targets.
A companion paper, \textit{On Disease Epistemology in Blind Receiver}, proves that
loop-holonomy self-consistency is the unique viable diagnostic criterion for
disease state identification (AUC$\,=1.00$, 25/25 validation experiments) and
provides a sparse $\ell_1$ linear program for therapeutic target selection with
minimum off-target effects; it is affiliated with the AIMe Registry and TUM
School of Life Sciences Weihenstephan.

\textbf{Cancer cell state transitions.}
The \href{https://github.com/fullscreen-triangle/syndrome}{Syndrome} framework
models cellular state through eight oscillator classes (protein, enzyme, membrane,
channel, ATP, genetic expression, calcium, circadian), each with a coherence index
$\eta_i \in [0,1]$, computing a disease vector $\mathbf{D} \in [0,1]^8$ that
locates the cell's current state in a partition geometry on the $[0,1]^8$ hypercube.
Trajectory computation runs in $\mathcal{O}(N\log N)$; therapeutic target selection
uses sparse LP to identify the minimum-intervention path from disease state to
healthy attractor. Applied to cancer biology, the framework addresses cell fate
transitions and metastatic progression as trajectory completion problems within
this state space --- a natural extension of the institute's interest in cell fate
transitions and the molecular mechanisms of metastasis.

\textbf{Multi-omics and imaging.}
For the metabolomics and imaging dimensions of multi-dimensional data integration,
\href{https://lavoisier.vercel.app/experiment}{lavoisier.vercel.app/experiment}
provides a virtual mass spectrometry instrument running force-free simulation
in the browser, with a downloadable Rust executable for large-scale local runs.
The underlying framework processes mzML datasets exceeding 100\,GB via Ray
distributed computing and reaches 98.9\,\% feature extraction accuracy against
the NIST library.
For protein-level analysis, \href{https://shakespear-nine.vercel.app/instrument}{shakespear-nine.vercel.app/instrument}
provides a browser-based instrument for protein folding, molecular trajectory,
catalysis, and dynamics via Kuramoto oscillator synchronization in S-entropy
space --- directly applicable to the oncoproteins and cancer-associated structural
changes identified from CRISPR screen readouts.
For small-molecule and drug target analysis,
\href{https://honjo-masamune.vercel.app/tools}{honjo-masamune.vercel.app/tools}
provides zero-parameter GPU-accelerated molecular spectroscopy including docking,
3D structure, fingerprints, and network analysis, running entirely in the browser.

\textbf{Institutional context.}
I completed doctoral research in Computational Lipidomics at TUM / Bayerisches
Forschungsinstitut, where I developed mass spectrometry pipelines and federated
learning architectures for distributed clinical analysis. Published work is
affiliated with the AIMe Registry and TUM School of Life Sciences Weihenstephan.
I am legally resident in Germany with full work authorisation and am available
immediately. English is native; German is conversational.

\vspace{1.5em}
Yours sincerely,

\vspace{2.5em}
\textbf{Kundai Farai Sachikonye}

\end{document}
